\begin{corollary}[6.3.2]
\label{IV.2.s6.3.2-fr}
Bajo las hipótesis de \textup{(6.3.1)}, supongamos además $M\neq0$ y $N\neq0$;
entonces se tiene
\begin{equation}
\operatorname{coprof}_B(M\otimes_AN)=\operatorname{coprof}_A(M)
+\operatorname{coprof}_{B\otimes_Ak}(N\otimes_Ak).
\tag{6.3.2.1}\label{IV.2.s6.3.2.1-fr}
\end{equation}
\end{corollary}

Esto resulta inmediatamente de (6.3.1.1), de (6.1.2) y de la definición de la
coprofundidad (\textbf 0, 16.4.9).

\begin{corollary}[6.3.3]
\label{IV.2.s6.3.3-fr}
Bajo las hipótesis de \textup{(6.3.1)}, supongamos además $M\neq0$ y $N\neq0$; para
que $M\otimes_AN$ sea un $B$-módulo de Cohen--Macaulay, es necesario y
suficiente que $M$ sea un $A$-módulo de Cohen--Macaulay y que $N\otimes_Ak$
sea un $(B\otimes_Ak)$-módulo de Cohen--Macaulay.
\end{corollary}

Esto resulta del corolario (6.3.2) y de la definición de los módulos de
Cohen--Macaulay por el hecho de que su coprofundidad es nula (\textbf 0, 16.5.1),
teniendo en cuenta que la condición $N\neq0$ equivale a $N\otimes_Ak\neq0$ en
virtud del lema de Nakayama.

\begin{corollary}[6.3.4]
\label{IV.2.s6.3.4-fr}
Bajo las hipótesis de \textup{(6.3.1)}, supongamos además que $N\otimes_Ak$ sea un
$B$-módulo de longitud finita; entonces se tiene
\begin{equation}
\operatorname{prof}_B(M\otimes_AN)=\operatorname{prof}_A(M).
\tag{6.3.4.1}\label{IV.2.s6.3.4.1-fr}
\end{equation}

\oldpage[IV-2]{140}

Si, además, $M\neq0$ y $N\neq0$, se tiene
\begin{equation}
\operatorname{coprof}_B(M\otimes_AN)=\operatorname{coprof}_A(M).
\tag{6.3.4.2}\label{IV.2.s6.3.4.2-fr}
\end{equation}
\end{corollary}

En efecto, equivale a lo mismo decir que $N\otimes_Ak$ es un $B$-módulo de
longitud finita o un $(B\otimes_Ak)$-módulo de longitud finita, y se sabe
(\textbf 0, 16.2.3) que los módulos de longitud finita $\neq0$ son de dimensión $0$ y
de profundidad $0$.

El corolario (6.3.4) será aplicable, en particular, cuando $B\otimes_Ak$ sea
un anillo de longitud finita, es decir, cuando $\mathfrak mB$ sea un \emph{ideal de
definición} del anillo $B$.
